\documentclass[11pt]{article}
\usepackage{amsmath,amssymb,amsthm,hyperref}
\newtheorem{lemma}{Lemma}
\title{Polynomial profile bridges between adaptive residue cones}
\author{}\date{}
\begin{document}\maketitle
\paragraph{Status.}
This note gives exact written identities and a finite rational certificate.
It does not establish an induction, a superconstant local saving, or either
accepted asymptotic conclusion. The model and its definitions are unchanged.
The certificates are independently checked rational arithmetic, not Lean proofs.

\paragraph{The finite bridge.}
Use the sixteen-pole grid and four target words in
\texttt{butterfly-biased-cone-intersection.json}, with
$p_y=\sum_{j=0}^3y_j\epsilon^{2^j}$ and $\epsilon=1/100$.
Let $K_0,\ldots,K_3$ be their open homogeneous residue cones after
eliminating the separate block biases, as in
\texttt{butterfly-adaptive-haar-biased.tex}.
Put $C_{ij}=(1-p_i-p_j)^{-1}$ and
\[
 V_2=C^{-1}\operatorname{span}\{(1)_i,(p_i)_i,(p_i^2)_i\}.
\]
For each of the three partitions $\{a,b\}\sqcup\{c,d\}$ of
$\{0,1,2,3\}$, the following intersection is nonempty:
\[
 (K_a\cap K_b)\cap\bigl((K_c\cap K_d)-V_2\bigr).             \tag{1}
\]
This removes the previously certified obstruction for the one-dimensional
space $C^{-1}\operatorname{span}\{(1)_i\}$ on these four targets.
It says nothing about all target words or all sufficiently small scales.

Here is a solver-free certificate of (1).
The file \texttt{butterfly-biased-polynomial-bridge.json} stores the four
words, the pole grid, the three partitions, and a square active basis
for each partition. For $V=C^{-1}[(1),(p),(p^2)]$, the unknowns are
$(x,\lambda)\in\mathbb Q^{16}\times\mathbb Q^3$.
A comparison row $r$ from the first two cones contributes $(r,0)$;
a row from the other two contributes $(r,rV)$.
Thus positivity of all rows means that $x$ is in the first intersection
and $x+V\lambda$ is in the second. The checker
\texttt{butterfly-biased-polynomial-bridge-check.py} solves the saved
active rows with right side one, together with the saved zero coordinate
constraints, and checks all original inequalities strictly. It checks
$264$ strict comparisons in total. No optimization solver or numerical
rounding is used by this checker.

\begin{lemma}[Inverse Cauchy polynomial spaces]
Let $p_1,\ldots,p_M$ be distinct, let $t_i=T-p_i$ be distinct and
outside the pole set, and let $C_{ij}=(t_i-p_j)^{-1}$.
Define
\[
 Q(x)=\prod_j(x-p_j),\qquad A(x)=\prod_i(x-t_i),\qquad
 w_j=-\frac{A(p_j)}{Q'(p_j)}.
\]
For $0\leq L<M$,
\[
 C^{-1}\operatorname{span}\{(p_i^k)_i:0\leq k\leq L\}
 =\operatorname{span}\{(w_jp_j^k)_j:0\leq k\leq L\}.          \tag{2}
\]
\end{lemma}
\begin{proof}
For a polynomial $P$ of degree at most $L$, divide
$Q(x)P(T-x)$ by the monic polynomial $A(x)$:
\[
 Q(x)P(T-x)=A(x)S(x)+R(x),\qquad \deg R<M,\quad\deg S\leq L.
\]
The rational function $R/Q$ takes the values $P(p_i)$ at $t_i$.
Its residue at $p_j$ is $-A(p_j)S(p_j)/Q'(p_j)=w_jS(p_j)$.
The quotient map from $P$ to $S$ is an invertible triangular map on
polynomials of degree at most $L$: its leading coefficient changes
only by the nonzero factor $(-1)^{\deg P}$. This proves both inclusions.
Invertibility of $C$ also follows directly: a null vector would give a
rational function with numerator degree less than $M$ and $M$ distinct
zeros outside the poles, hence zero residues.
\end{proof}

\paragraph{Binary clustering does not preserve these weights directly.}
Replace each child pole $p_j$ by $p_j,p_j+\delta$, keeping the reflected
query convention. The parent polynomials are
\[
 Q_\delta(x)=Q(x)Q(x-\delta),\qquad
 A_\delta(x)=A(x)A(x+\delta).
\]
For fixed distinct child poles and $\delta\to0^+$, their weights obey
\[
 w_{j,0}(\delta)=\frac{w_j^2}{\delta}+O(1),\qquad
 w_{j,1}(\delta)=-\frac{w_j^2}{\delta}+O(1).                  \tag{3}
\]
Indeed $Q(p_j-\delta)=-\delta Q'(p_j)+O(\delta^2)$ and
$Q(p_j+\delta)=\delta Q'(p_j)+O(\delta^2)$; substitution into the
weight formula gives (3).
Thus the most direct pair-difference limit gives squared child weights.
Converting these into the weights in (2) requires the sampled factor
$1/w_j=-Q'(p_j)/A(p_j)$, which has no established bounded-degree
polynomial representation on the child grid. This is a gap in the
proposed induction, not an impossibility theorem for other profiles.

\begin{lemma}[Two heads for an arbitrary fixed query profile]
For any row values $b_1,\ldots,b_r\in\mathbb R$ and fixed query profile
$v_1,\ldots,v_s\in\mathbb R$, there is an axis-zero sum of two
strictly admissible reverse heads taking the value $b_iv_j$ at
$(e_i,e_j)$.
\end{lemma}
\begin{proof}
Put $C=2r+2s+4$. Choose $\kappa>0$ small enough that
\[
 u_j=C-1-\left(\frac1{C-2}+\kappa v_j\right)^{-1}
 \quad\hbox{satisfies}\quad \tfrac12<u_j<\tfrac32
\]
for every $j$. This is possible by continuity because $u_j=1$ at
$\kappa=0$. Set
\[
 F(z,w)=\frac{\kappa^{-1}\sum_i b_iz_i}{C-\sum_i z_i-\sum_j u_jw_j}
       -\frac{\kappa^{-1}\sum_i b_iz_i}{C-\sum_i z_i-\sum_j w_j}.
\]
Both denominators have strictly negative slopes and are positive on the
entire Boolean cube, since their minimum is at least $C-r-3s/2>0$.
At $(e_i,e_j)$ the reciprocal denominator difference is exactly
$\kappa v_j$, so $F=b_iv_j$. At $z=0$ the numerators vanish, and at
$w=0$ the denominators agree. Both Boolean coordinate axes therefore
vanish exactly.
\end{proof}

\paragraph{Budget and remaining issue.}
If a block at butterfly level $k$ is allowed $B_k$ fixed query profiles,
the preceding lemma implements all such profile corrections with
$2\sum_{k=0}^{d-1}H B_k/2^{k+1}$ heads. This is $O(H)$ for
$B_k=O(k+1)$ and also for $B_k=O(2^{k/2})$.
Thus these profile dimensions would preserve a hypothetical
$H\log H$ decoder saving. What is missing is a universal cone-intersection
or sign-realization theorem with any such profile budget. The finite
certificate (1) and identity (2) do not supply that theorem.
\end{document}
